% Part: lambda-calculus
% Chapter: introduction
% Section: overview

\documentclass[../../../include/open-logic-section]{subfiles}

\begin{document}

\olfileid{lam}{int}{ovr}
\olsection{نظرة عامة}

وضع ألونزو تشيرش حساب لامبدا في أوائل ثلاثينيات القرن العشرين
ليكون أساسًا للمنطق البنائي، \emph{لا} ليكون نموذجًا للدوال القابلة
للحساب. ثم لم يلبث أن ثبت تكافؤه مع تعريفات أخرى للقابلية للحساب،
منها دوال تورنغ القابلة للحساب والدوال العودية الجزئية.
وكان في هذا التكافؤ يومئذ شيء من المفاجأة، فزاد توصيف القابلية
للحساب به جدارة بالنظر.

بترميز لامبدا نشير إلى الدالة بالتعبير الرمزي الذي يعرّفها نفسه،
فلا نحتاج أولًا إلى تسميتها باسم مستقل. فبدل قولنا:
«لتكن~$f$ الدالة المعرفة بالمعادلة~$f(x) = x + 3$»، نقول:
«لتكن~$f$ الدالة~$\lambd[x][(x + 3)]$».
فالتعبير~$\lambd[x][(x+3)]$ \emph{اسم} للدالة التي تزيد وسيطها
ثلاثة. وليس~$x$ فيه إلا متغيرًا صوريًا يشغل موضع الوسيط؛
ولذلك يصح أن ندل على الدالة نفسها بـ$\lambd[y][(y + 3)]$.
ولا يبطل هذا الترميز بوجود معلمات أخرى.
فإذا كانت~$g(x, y)$ دالة في متغيرين، وكان~$k$ عددًا طبيعيًا،
كانت~$\lambd[x][g(x,k)]$ الدالة التي ترسل كل~$x$ إلى~$g(x, k)$.

هذه الطريقة في إنشاء اسم الدالة من تعبير رمزي هي
\emph{تجريد لامبدا}. ويقابل التجريد \emph{التطبيق}:
متى أعطينا دالة~$f$، على الأعداد الطبيعية مثلًا، أمكن
\emph{تطبيقها} على قيمة مثل 2، فنكتب النتيجة في الترميز
المعتاد~$f(2)$.

وإذا جمعنا التجريد والتطبيق أمكن تبسيط الناتج بإحلال الوسيط
المطبّق عليه محل المتغير المجرّد. فالتعبير
\[
(\lambd[x][(x + 3)])(2)
\]
يُبسّط إلى~$2 + 3$.

إلى هنا لم نزد على مفاهيم مألوفة إلا ترميزًا جديدًا.
لكن حساب لامبدا نفسه يفارق تصور نظرية المجموعات مفارقة أعمق؛
ففيه:
\begin{enumerate}
\item كل تعبير يدل على دالة.
\item تُعرّف الدوال بتجريد لامبدا.
\item يجوز تطبيق أي شيء على أي شيء آخر.
\end{enumerate}
فمتى كان~$F$ حدًا من حدوده، عُدّ~$F(F)$ ذا معنى دائمًا.
وهذا الإطار الذي يجيز التطبيق بلا تلك القيود يسمى حساب
لامبدا \emph{غير المنمّط}؛ ومعنى ذلك أنه لا قيد فيه على ما
يطبّق على ماذا.

\begin{digress}
ولحساب لامبدا صورة أخرى مهمة هي الحساب \emph{المنمّط}.
تشبه الصورة غير المنمطة من وجوه كثيرة، لكنها أيسر كثيرًا في
التوفيق مع الإطار الكلاسيكي لنظرية المجموعات، وتختلف عنها في
بعض الخصائص اختلافًا بعيدًا.

وقد صار البحث في حساب لامبدا من أسس علم الحاسوب النظري وتصميم
لغات البرمجة. ولغة LISP، التي وضعها جون مكارثي في خمسينيات
القرن العشرين، مثال مبكر للغات التي أخذت بهذه الأفكار.
\end{digress}

\end{document}
